Um zu zeigen, dass $(\sigma \cdot \mathbf{p})^2 = \mathbf{p}^2$, betrachten wir die Pauli-Matrizen $\sigma_i$ und den Impulsoperator $\mathbf{p} = (p_1, p_2, p_3)$. Das Skalarprodukt $\sigma \cdot \mathbf{p}$ ist definiert als $\sigma \cdot \mathbf{p} = \sigma_1 p_1 + \sigma_2 p_2 + \sigma_3 p_3$. Wir berechnen nun das Quadrat:

\[
(\sigma \cdot \mathbf{p})^2 = (\sigma_1 p_1 + \sigma_2 p_2 + \sigma_3 p_3)^2.
\]

Unter Verwendung der Eigenschaften der Pauli-Matrizen, insbesondere $\sigma_i \sigma_j = \delta_{ij} + i \epsilon_{ijk} \sigma_k$, ergibt sich:

\[
(\sigma \cdot \mathbf{p})^2 = \sum_{i=1}^{3} \sigma_i^2 p_i^2 + \sum_{i \neq j} \sigma_i \sigma_j p_i p_j.
\]

Da $\sigma_i^2 = I$ (die Einheitsmatrix) und $\sigma_i \sigma_j + \sigma_j \sigma_i = 0$ für $i \neq j$, folgt:

\[
(\sigma \cdot \mathbf{p})^2 = \sum_{i=1}^{3} p_i^2 = \mathbf{p}^2.
\]

Für die Spinoren der Dirac-Gleichung zeigen wir, dass $u_r^{+}(p) u_s(p) = v_r^{+}(p) v_s(p) = 2 E \delta_{rs}$. Die Dirac-Spinoren $u_r(p)$ und $v_r(p)$ sind Lösungen der Dirac-Gleichung. Die Normierungsbedingung ergibt sich aus der Forderung, dass die Spinoren orthonormal sind und die Energie $E$ berücksichtigen:

\[
u_r^{+}(p) u_s(p) = v_r^{+}(p) v_s(p) = 2 E \delta_{rs}.
\]

Dies folgt aus der kanonischen Quantisierung und der Erhaltung der Wahrscheinlichkeit.

Um zu zeigen, dass $\left\{\gamma^\mu, \gamma^\nu\right\} = 2 g^{\mu \nu}$, verwenden wir die Definition der Dirac-Matrizen $\gamma^\mu$ und die Antikommutator-Relation:

\[
\left\{\gamma^\mu, \gamma^\nu\right\} = \gamma^\mu \gamma^\nu + \gamma^\nu \gamma^\mu.
\]

Die Dirac-Matrizen sind so konstruiert, dass sie die Clifford-Algebra erfüllen, was bedeutet, dass:

\[
\gamma^\mu \gamma^\nu + \gamma^\nu \gamma^\mu = 2 g^{\mu \nu} I,
\]

wobei $g^{\mu \nu}$ der metrische Tensor der Raumzeit ist. Dies stellt sicher, dass die Dirac-Matrizen die korrekten Transformationseigenschaften unter Lorentz-Transformationen haben.